\documentclass[11pt,a4paper]{article}
\usepackage[T1]{fontenc}
\usepackage{lmodern,amsmath,amssymb}
\usepackage[top=18mm,bottom=24mm,left=22mm,right=22mm,headheight=14pt,footskip=12mm]{geometry}
\usepackage{microtype}
\usepackage{array,booktabs,longtable}
\usepackage[table]{xcolor}
\usepackage{parskip,titlesec,fancyhdr}
\usepackage{hyperref}
\input{glyphtounicode}
\pdfgentounicode=1
\definecolor{ink}{HTML}{192A38}
\definecolor{teal}{HTML}{145F58}
\definecolor{ruleteal}{HTML}{B8CFCD}
\definecolor{wash}{HTML}{EAF4F1}
\hypersetup{colorlinks,urlcolor=teal,pdftitle={Induced-universal posets of height three},pdfsubject={The Mantel–Hall construction}}
\setlength{\parskip}{5pt plus 1pt}
\setlength{\parindent}{0pt}
\linespread{1.07}
\setcounter{secnumdepth}{0}
\titleformat{\section}{\sffamily\Large\bfseries\color{teal}}{}{0pt}{}
\titleformat{\subsection}{\sffamily\normalsize\bfseries\color{teal}}{}{0pt}{}
\titlespacing*{\section}{0pt}{13pt plus 2pt}{5pt}
\titlespacing*{\subsection}{0pt}{9pt}{2pt}
\pagestyle{fancy}
\fancyhf{}
\renewcommand{\headrulewidth}{0pt}
\renewcommand{\footrulewidth}{0.4pt}
\makeatletter
\renewcommand{\footrule}{\hbox to\headwidth{\color{ruleteal}\leaders\hrule height \footrulewidth\hfill}\vskip 4pt}
\makeatother
\fancyfoot[L]{\sffamily\footnotesize\color{teal}UNIVERSAL POSETS\enspace /\enspace HEIGHT THREE}
\fancyfoot[R]{\sffamily\footnotesize\color{teal}\thepage}
\newcommand{\tightlist}{\setlength{\itemsep}{2pt}\setlength{\parskip}{0pt}}
\newcommand{\PaperTitle}{%
  {\sffamily\footnotesize\bfseries\color{teal}MANTEL--HALL CONSTRUCTION\par}
  \vspace{6pt}
  {\sffamily\bfseries\fontsize{30}{32}\selectfont Induced-universal posets\\of height three\par}
  \vspace{12pt}}
\newenvironment{resultbox}{\par\begingroup\setlength{\fboxsep}{11pt}\noindent
  \begin{lrbox}{0}\begin{minipage}{\dimexpr\linewidth-2\fboxsep\relax}\setlength{\parskip}{4pt}}
  {\end{minipage}\end{lrbox}\colorbox{wash}{\usebox{0}}\par\endgroup\vspace{3pt}}
\begin{document}
\color{ink}
\PaperTitle

\begin{resultbox}

\textbf{Theorem.} Let \(u_3(n)\) be the least size of a poset of height
at most three that contains every \(n\)-element poset of height at most
three as an induced subposet. Then \[
u_3(n)=2^{n/4+o(n)}.
\] The same exponent holds when the targets are restricted to connected
posets whose maximal chains all have three elements.

\end{resultbox}

\section{1. Encode the order as a triangle-free
graph}\label{encode-the-order-as-a-triangle-free-graph}

Assign each element the length of the longest chain ending there. Denote
the three ranks by \(A,M,C\) from bottom to top. Empty ranks are
allowed. Every strict comparison goes up in rank. An induced embedding
must preserve both comparisons and incomparabilities.

\begin{longtable}[]{@{}ll@{}}
\toprule\noalign{}
Rank pair & A graph edge means \\
\midrule\noalign{}
\endhead
\bottomrule\noalign{}
\endlastfoot
Bottom--middle, \(A,M\) & \(a<m\) \\
Middle--top, \(M,C\) & \(m<c\) \\
Bottom--top, \(A,C\) & \(a\) and \(c\) are incomparable \\
\end{longtable}

A triangle in the resulting graph \(G\) would give \(a<m<c\) with
incomparable outer endpoints. Conversely, every triangle-free graph on
these three ranks defines a poset by the same rules. The only two-step
chains to check traverse all three ranks, and triangle-freeness forces
the composite comparison. The rule also allows a bottom--top comparison
with no middle element between its endpoints.

\section{2. Assign bits to slots}\label{assign-bits-to-slots}

A \textbf{slot} is a named position on a rank. Fix a triangle-free
tripartite graph \(H\) on at most \(n\) slots. For each cross-rank pair
in \(H\), designate one endpoint to read its graph bit. For every other
cross-rank pair, designate both endpoints. We call \(H\) the template.

Give each slot \(b=\lceil n/4\rceil\) copies. Match each edge of \(H\)
to a copy of either endpoint. Any set of edges with endpoint set \(W\)
has at most \(e(H[W])\) edges. Mantel's theorem yields \[
e(H[W])\le |W|^2/4\le b|W|.
\] The neighborhood in the matching problem has exactly \(b|W|\) copies,
so Hall's theorem gives an assignment. Fix it for the template. Each
slot reads at most \(b\) bits from pairs in \(H\).

\newpage

\section{3. Decode comparisons on arbitrary
labels}\label{decode-comparisons-on-arbitrary-labels}

A host vertex is a label at a slot. For a label \(x\), let \(S_x\) be
the set of slots whose bits it reads as \(1\). In expressions involving
\(S_x\), a label name can also stand for its slot. The host order goes
strictly upward in rank.

\begin{itemize}
\tightlist
\item
  \textbf{Adjacent ranks.} Compare the lower label to the upper label if
  all designated readings of the pair are \(1\) and the labels have no
  common positive slot on the third rank.
\item
  \textbf{Bottom to top.} Declare \(a<c\) if any designated reading of
  \(ac\) is \(0\) or if \(S_a\cap S_c\) contains a middle slot.
\end{itemize}

\subsection{Transitivity}\label{transitivity}

Suppose three arbitrary labels satisfy \(a<m<c\). If any reading of
\(ac\) is \(0\), the outer rule gives \(a<c\). Otherwise, choose an
outer endpoint \(x\) that reads \(ac\) as \(1\) and call the other
endpoint \(y\). Thus \(y\in S_x\).

If \(m\) read \(my\), the comparison between \(m\) and \(y\) would force
its reading to be \(1\). But then \(y\in S_x\cap S_m\) would veto the
comparison between \(x\) and \(m\). Therefore \(m\) does not read
\(my\). This forces \(my\in H\), with \(y\) as its sole reader, and
\(m\in S_y\).

\begin{itemize}
\tightlist
\item
  If \(ac\in H\), the two edges \(ac,my\in H\) force \(mx\notin H\) by
  triangle-freeness. The slot \(x\) reads \(mx\), and the adjacent
  comparison forces that reading to be \(1\).
\item
  If \(ac\notin H\), both outer endpoints read \(ac\) as \(1\). Applying
  the same veto argument with \(x\) and \(y\) exchanged forces \(x\) to
  read \(mx\) as \(1\).
\end{itemize}

Thus \(m\in S_a\cap S_c\), and the outer rule gives \(a<c\). This proof
applies to every combination of labels, even if two labels disagree on a
pair read twice. Strict rank increase gives irreflexivity and height at
most three.

\subsection{Induced embeddings and label
count}\label{induced-embeddings-and-label-count}

To embed a target, select the label at each occupied slot that records
its genuine bits. An adjacent edge passes the reading test and cannot be
vetoed, since a common positive slot would give a triangle in \(G\). A
missing adjacent edge fails the reading test. An outer comparison has
bit \(0\). An outer incomparability has bit \(1\), and a common positive
middle slot would again yield a triangle. The decoder therefore recovers
the whole order, and distinct slots give an injection.

Suppose each slot has at most \(d\) positive bits outside \(H\). Its
label records the assigned bits on \(H\) explicitly and lists the
positions of its positive bits elsewhere. Including every label allowed
by these budgets gives at most \[
n\,2^{\lceil n/4\rceil}\sum_{i=0}^{d}\binom ni
\] vertices for one template.

\newpage

\section{4. Find a bounded family of
templates}\label{find-a-bounded-family-of-templates}

Fix \(0<\delta<1/64\) and \(0<\varepsilon<\delta/4\). The rank-refining
equitable Szemerédi regularity lemma {[}1{]} gives, for sufficiently
large \(n\), an exceptional set \(V_0\) of size at most
\(\varepsilon n\) and \(t\le T(\varepsilon)\) clusters of equal size
\(q\). Each cluster lies in one rank and at most \(\varepsilon t^2\)
cluster pairs are irregular. The bound \(T(\varepsilon)\) is independent
of \(n\).

Join two clusters in the reduced graph when their pair is
\(\varepsilon\)-regular and has density at least \(\delta\). The reduced
graph is triangle-free. In a putative triangle of clusters \(X,Y,Z\),
fewer than \(\varepsilon q\) vertices in \(X\) have fewer than
\((\delta-\varepsilon)q\) neighbors in \(Y\), and the same bound holds
for \(Z\). Choose \(x\) outside both bad sets. The two neighbor sets
each have size at least \(\varepsilon q\). Regularity between \(Y,Z\)
gives density at least \(\delta-\varepsilon>0\) between those sets, so
an edge between them makes a triangle in \(G\).

Complete every joined cluster pair to a biclique to form \(H\). There
are no template edges incident to \(V_0\). A triangle in \(H\) would
project to a triangle of clusters, so \(H\) is triangle-free.

\subsection{Remove high-degree
deviations}\label{remove-high-degree-deviations}

The edges of \(G\) outside \(H\) are the positive bits missing from the
template. The low-density regular pairs contribute at most
\(\delta n^2\) such edges, irregular pairs at most
\(\varepsilon t^2q^2\le\varepsilon n^2\), and pairs touching \(V_0\) at
most \(\varepsilon n^2\). Their total is at most \(2\delta n^2\).

By summing degrees, at most \(4\sqrt\delta\,n\) vertices have more than
\(\sqrt\delta\,n\) uncovered neighbors. Set those vertices and \(V_0\)
aside. There are at most \[
k=\left\lfloor(4\sqrt\delta+\varepsilon)n\right\rfloor
\] exceptional vertices, and every surviving slot has at most
\(d=\lfloor\sqrt\delta\,n\rfloor\) positive bits outside \(H\).

\subsection{Use a fixed number of slots per
cluster}\label{use-a-fixed-number-of-slots-per-cluster}

Give every cluster \(\lfloor n/t\rfloor\) slots. This is at least \(q\),
so survivors fit and all remaining slots can be left vacant. Set every
target graph bit touching a vacant slot to \(0\). These slots are
isolated in the extended target graph, which remains triangle-free, and
they add no positive bits outside \(H\). The template uses at most \(n\)
slots in total.

A template now depends only on \(t\), the rank of each cluster, and the
reduced graph. Thus there are at most \[
C_T=\sum_{t=1}^{T(\varepsilon)}3^t2^{\binom t2}
\] templates. This is a constant for fixed \(\varepsilon\). Take their
disjoint union as the core host \(B\). Its size is at most \(C_T\) times
the label count from Section 3.

\newpage

\section{5. Restore the exceptional
vertices}\label{restore-the-exceptional-vertices}

Use a single pool \([k]\) of exception identifiers. Decorate each core
vertex with every pair of disjoint sets \(D,U\subseteq[k]\). In an
embedding they record the exceptional vertices below and above it. For
each identifier \(i\), include a label at each of the three ranks for
every pair with \(D\cap U=\{i\}\). An exceptional vertex counts itself
in both sets. Unused identifiers are absent from the genuine signatures.

For labels \(x,y\) on strictly increasing ranks, declare \(x<y\) if \[
D_x\subseteq D_y,\qquad U_y\subseteq U_x,
\] and either their core vertices compare in \(B\), or
\(U_x\cap D_y\ne\varnothing\). The core clause applies only when both
labels have core vertices.

\textbf{Transitivity.} If \(x<y<z\), the inclusions compose. A witness
in \(U_x\cap D_y\) survives because \(D_y\subseteq D_z\). A witness in
\(U_y\cap D_z\) survives because \(U_y\subseteq U_x\). If neither step
uses a witness, both steps come from \(B\), whose transitivity gives the
composite. The rank condition still bounds the height by three.

\textbf{Induced embedding.} Use the core embedding and the actual
signatures and ranks. A true comparison yields the inclusions. If both
endpoints are core vertices, their core comparison suffices. If one
endpoint is exceptional, it witnesses the comparison itself. Conversely,
a declared comparison either comes from the induced core embedding or
has a genuine exception weakly between the endpoints. Either way it is a
true comparison.

For a core vertex there are \(3^k\) signatures. For each exception
identifier there are three ranks and \(3^{k-1}\) signatures per rank
when \(k\ge1\). Thus the extension has at most \(3^k(|B|+k)\) vertices,
with the same formula at \(k=0\).

\section{6. Finish the upper bound}\label{finish-the-upper-bound}

For \(\rho=\sqrt\delta<1/2\), the binomial sum is at most
\(2^{h_2(\rho)n}\), where
\(h_2(\rho)=-\rho\log_2\rho-(1-\rho)\log_2(1-\rho)\). Indeed,
multiplying by \(s^{\rho n}\) for \(s=\rho/(1-\rho)<1\) bounds the
resulting sum by \((1+s)^n\). Combining all factors gives \[
\log_2 u_3(n)\le
\frac n4+
\bigl[h_2(\sqrt\delta)+(4\sqrt\delta+\varepsilon)\log_2 3\bigr]n
+\log_2 n+O_{\delta,\varepsilon}(1).
\] For any prescribed tolerance, choose \(\delta\) and \(\varepsilon\)
so that the bracket is smaller than half that tolerance. Then take \(n\)
large enough to absorb the logarithm and the constant. This proves
\(u_3(n)\le 2^{n/4+o(n)}\).

\newpage

\section{7. Match the exponent by
counting}\label{match-the-exponent-by-counting}

Use one bottom element \(a\), and split the other \(n-1\) labeled
elements as evenly as possible between \(M\) and \(C\). Put \(a\) below
every other element. In the \(M\times C\) comparison matrix, fix one row
and one column to all \(1\)s. Choose all other entries independently.

The fixed row ensures that each top element has a middle predecessor.
The fixed column gives every middle element a top successor. With the
single bottom element, this makes every maximal chain have three
elements and makes the order connected.

The free entries give \[
2^{(|M|-1)(|C|-1)}=2^{n^2/4-O(n)}
\] different labeled orders. A host on \(N\) vertices admits at most
\(N^n\) injections from the \(n\) labels, and each injection induces at
most one order. Thus \[
N^n\ge 2^{n^2/4-O(n)}
\quad\Longrightarrow\quad
N\ge 2^{n/4-O(1)}.
\] Combining this with the upper bound proves the theorem for all
height-three targets and for the connected pure subfamily. \(\square\)

\subsection{Proof and accompanying
material}\label{proof-and-accompanying-material}

The construction gives the optimal leading exponent for the full family.
The proof uses Szemerédi regularity, so its overhead is \(2^{o(n)}\).
The repository note separately treats the bipartite signed-graph family,
sharper exceptional-vertex pooling, and the 657-vertex benchmark for the
rank profile \((4,8,4)\).

The Lean source states the asymptotic theorem and includes the decoder,
extension, and counting declarations. The integration of the regularity
argument into the formal asymptotic proof remains open.

\begin{itemize}
\tightlist
\item
  \href{https://github.com/hyper-analysis/universal-poset-mantel/blob/main/math/height-three-universal-posets.md}{Repository
  proof notes and refinements}
\item
  \href{https://github.com/hyper-analysis/universal-poset-mantel/blob/main/math/HeightThreeBound.lean}{Lean
  source}
\item
  \href{https://github.com/hyper-analysis/universal-poset-mantel/tree/main/experiments/height-three-universal-posets/profile-benchmark}{Fixed-profile
  benchmark and replay}
\end{itemize}

\subsection{References}\label{references}

{[}1{]} József Balogh, Ramon I. Garcia, and Marcelo Sales,
\href{https://arxiv.org/abs/2607.12980}{\emph{Even smaller universal
posets}}, arXiv:2607.12980. Theorem 3.1 gives the refining regularity
lemma used here.
\end{document}
